\clearpage
\fancyhead[L]{\small CERES PALLAS.}
\SourceLeaf{467}
\begin{center}
\rule{0.55\linewidth}{0.4pt}\\[2pt]
Astronomische Nachrichten Band VII. Nr. 161. Seite 329. 1829 Juni.\\[2pt]
\rule{0.55\linewidth}{0.4pt}
\end{center}

\begin{center}[Beobachtungen der \emph{Ceres} in Göttingen von \textsc{Gauss.}]\end{center}

\begin{center}
\begin{tabular}{l|l|l|l}
1829 & Mittlere Zeit & Rektaszension & Nördliche Deklination \\
\hline
Januar 22 & $10^{\rm u}21^{\rm m}46^{\rm s}2$ & $97^\circ24'25''0$ & $30^\circ22'0''4$ \\
\hfill 25 & $10\;7\;41.4$ & $96\;50\;4.0$ & $30\;30\;8.3$ \\
Februar 11 & $8\;52\;51.7$ & $94\;49\;55.8$ & $30\;58\;55.4$ \\
\hfill 19 & $8\;20\;52.8$ & $94\;41\;58.8$ & $31\;4\;46.7$ \\
März 5 & $7\;29\;47.9$ & $95\;41\;35.9$ & $31\;6\;19.2$ \\
\hfill 14 & $6\;59\;57.5$ & $97\;5\;4.0$ & $31\;2\;42.0$ \\
\hfill 16 & $6\;53\;37.2$ & $97\;27\;56.7$ & \\
\hfill 18 & $6\;47\;22.3$ & $97\;52\;14.5$ & $30\;59\;55.9$ \\
\end{tabular}
\end{center}

Anhaltend ungünstiges Wetter verhinderte die Beobachtung der Opposition.

